\chapter{INTRODUCTION}

% ================================================
\section{Motivation}
% ================================================

Modern data-intensive systems rely on real-time data streams for decision-making, operational control, and downstream analytics. When these streams contain erroneous, incomplete, or anomalous records, the damage propagates rapidly through every dependent system in a \textbf{contamination cascade}: a single erroneous transaction contaminates payment processing, financial reporting, and fraud detection models simultaneously. Industry analysis documents eight distinct failure modes in streaming data transmission: corrupted data from encoding errors, schemaless events lacking structural metadata, invalid schemas, incompatible schema evolution, logically invalid field values, semantically incorrect values, missing events, and duplicate events \citep{confluent_2025}. The financial stakes of poor data quality are substantial and quantifiable. According to IBM's 2024 Data Quality Report, organizations lose an average of USD 12.9 million annually due to direct remediation costs, analyst productivity losses, poor decision-making from erroneous data, and regulatory penalties \citep{ibm_dq2024}. IBM's cost breakdown identifies four principal categories: (a) \textbf{Direct remediation costs} for finding, diagnosing, and fixing data errors, including manual review labor and automated validation systems; (b) \textbf{Lost productivity}, where analysts and data engineers spend an estimated 40\% of their time on data preparation and cleaning rather than value-adding analysis; (c) \textbf{Poor decision-making}, including missed opportunities, suboptimal resource allocation, and customer experience degradation from decisions made on erroneous data; and (d) \textbf{Regulatory penalties}, including GDPR fines of up to 4\% of global annual revenue, SOX audit requirements, and industry-specific mandates. For financial services specifically, a single bad data incident at a major financial institution can result in USD 100 million or more in remediation costs, SEC enforcement actions, and potential civil litigation. The USD 4.7 billion figure from industry reports reflects cumulative enforcement actions across banking, healthcare, and telecommunications sectors by regulatory bodies including the SEC, OCC, FCA, and EU DPAs for systematic data quality failures --- representing averages across industries, with financial services and healthcare facing the highest per-incident costs due to regulatory scrutiny and downstream market integrity impacts. The data quality cost challenge intensifies as organizations adopt real-time analytics and AI-driven decision-making, because data quality issues propagate more rapidly through streaming pipelines where bad data at 10,000 events per second generates 600 errors per minute with no human review opportunity. An ounce of prevention is worth a pound of cure: catching bad data early prevents contamination from spreading to downstream systems (payment processing, financial reporting, fraud detection models), making real-time quality monitoring essential rather than optional.

The shift from batch processing to streaming paradigms fundamentally changes the nature of the data quality problem. In batch processing, all records are available before validation, enabling complete statistical analysis. In streaming, records arrive continuously and must be validated with minimal latency, requiring algorithms that update incrementally and make decisions under uncertainty with incomplete information. This epistemological shift renders traditional batch-oriented DQ frameworks architecturally incompatible with streaming contexts: they require complete datasets before validation, cannot operate on individual records as they arrive, and provide no mechanisms for online adaptation to distribution shifts.

Three forces converge to make streaming DQ monitoring both urgent and unsolved. First, the proliferation of high-throughput streaming platforms (Apache Kafka, Apache Flink, Confluent, Grab's Coban) means that more data is in motion than ever before, creating more attack surface for data quality failures. Second, the integration of machine learning into operational workflows means that downstream models consume streaming data directly, making ML performance contingent on streaming data quality. Third, the scale of modern deployments means that data quality issues that would be manageable in batch settings become catastrophic in streaming settings: an erroneous transaction at 10,000 events per second generates 600 errors per minute with no human review opportunity.

Despite this urgency, existing open-source DQ frameworks (Great Expectations, AWS Deequ, Soda Core) operate exclusively in batch mode, focusing on schema validation and constraint checking against static or periodically refreshed datasets. Recent extensions to streaming contexts (Stream DaQ, Grab Coban, Confluent) demonstrate feasibility but share a critical limitation: they rely primarily on deterministic rule-based validation without ML-based anomaly detection for multivariate violations, lack context-aware thresholding mechanisms calibrated to the specific conditions of each observation, and provide no systematic approach to concept drift adaptation.

% ================================================
\section{Research Challenges}
% ================================================

Addressing streaming DQ monitoring with ML-based, context-aware, and drift-adaptive capabilities confronts three fundamental research challenges that existing literature has not simultaneously resolved.

\textbf{Challenge 1: Context Collapse in Heterogeneous Data Streams.} Streaming data streams are inherently heterogeneous: financial transactions vary by account type and market conditions, healthcare monitoring varies by patient demographics and clinical context, and urban transportation data varies by trip type, time of day, day of week, and geographic zone. In all these domains, applying a single global threshold across heterogeneous observations causes context collapse: legitimate high-value events are misclassified as anomalies because the threshold was calibrated on the average population. The literature establishes that context-aware thresholding outperforms global thresholding in principle (DTD proves no fixed threshold is universally optimal \citep{dtd2026}; ReAD shows relative divergence metrics outperform fixed thresholds \citep{read2020}; BeSTAD demonstrates multi-dimensional behavioral profiles significantly outperform single-score approaches \citep{bestad2025}), yet no existing streaming DQ framework provides automated multi-dimensional context-aware threshold computation with drift-aware recalibration. The open question is how many contextual dimensions are necessary and sufficient, how to handle sparse context cells with insufficient training data, and how context-aware thresholds behave under concept drift.

\textbf{Challenge 2: Architectural Inefficiency in Validation Pipelines.} Existing streaming DQ frameworks process records through validation stages sequentially: Schema Validation $\rightarrow$ Business Rules $\rightarrow$ ML Scoring. Every record pays the full computational cost of the slowest stage regardless of whether earlier stages could make a determination. The key insight is that not all records require the same validation depth: most records fail early lightweight checks (schema violations, business rule breaches) and can be cheaply routed to the output stream without ever entering the expensive ML scoring stage. The open question is how to design a sequential pipeline that exploits early exit opportunities at lightweight stages to bypass expensive ML scoring, and under what violation rates this staged approach provides meaningful latency reduction.

\textbf{Challenge 3: Single-Strategy Drift Adaptation for Multiple Drift Types.} Production streaming systems encounter diverse drift types: gradual seasonal shifts, sudden event-driven spikes, and recurring cyclical patterns. Existing concept drift literature provides multiple adaptation strategies (periodic retraining, incremental updates, model switching, threshold adjustment) but applies them uniformly. This wastes resources on minor shifts that could be resolved by threshold adjustment while potentially under-responding to severe shifts that require full retraining. The literature establishes that automated strategy selection based on drift characteristics is theoretically promising (DTD proves dynamic thresholds strictly outperform fixed thresholds \citep{dtd2026}; ADWIN-U demonstrates unsupervised drift detection \citep{adwinu2025}; Adaptive-Delta ADWIN shows adaptive sensitivity reduces false alarms \citep{adt2025}) yet no existing framework combines drift detection with automated strategy selection.

These three challenges interact: context-aware thresholds computed from training data become stale under concept drift (Challenge 3), requiring drift monitoring; drift monitoring must operate within the throughput constraints of the processing architecture (Challenge 2); and the validation pipeline must provide enough throughput headroom to support context-aware ML scoring without blocking (Challenge 1). A unified solution must address all three simultaneously.

% ================================================
\section{Objective}
% ================================================

The primary objective of this thesis is to design, implement, and evaluate \textbf{CA-DQStream}, a \textbf{Context-Aware} framework for \textbf{Streaming Data Quality} monitoring that simultaneously addresses all three research challenges. Specifically, CA-DQStream aims to:

\begin{enumerate}
    \item Develop a \textbf{multi-dimensional context-aware thresholding framework} that partitions heterogeneous streaming data into context cells, computes individualized anomaly thresholds per cell automatically from training data, and adapts thresholds when data distributions drift. The goal is to reduce false positive rates compared to global thresholds while maintaining detection sensitivity across all context combinations.

    \item Design a \textbf{sequential pipeline architecture with staged validation} that exploits early exit opportunities at lightweight stages (schema, business rules) to bypass expensive ML scoring, reducing average processing latency and increasing throughput compared to naive sequential pipelines. The goal is to demonstrate that staged validation with early exit routing provides measurable performance benefits under realistic violation rates.

    \item Implement a \textbf{multi-strategy self-adaptation controller} that monitors quality meta-streams via drift detection and selects from multiple adaptation strategies automatically based on observed drift characteristics. The goal is to maintain detection accuracy under concept drift while minimizing computational cost by preferring lightweight strategies.

    \item Evaluate all innovations empirically on a \textbf{large-scale real-world streaming dataset}, measuring false positive rate, detection accuracy, throughput, latency, and drift adaptation effectiveness against established baselines and ablation configurations.
\end{enumerate}

% ================================================
\section{Contribution Highlights}
% ================================================

CA-DQStream makes four primary contributions, each addressing one of the research challenges and one or more of the research questions (RQ1--RQ5) stated in Chapter 2:

\begin{enumerate}
    \item \textbf{4D Context-Aware Thresholding Framework} [RQ1, RQ5]: A multi-dimensional threshold matrix partitioning the data space by trip type, time window, day type, and neighborhood with automated per-cell threshold computation from training data and intelligent fallback for sparse cells. Significantly reduces false positive rates compared to global thresholds, particularly for high-value data segments where global thresholds fail.

    \item \textbf{Sequential Four-Stage Processing Pipeline} [RQ2]: A sequential processing pipeline where lightweight validation stages (schema, business rules) filter out obvious violations early, allowing expensive ML scoring to operate only on records that pass prior stages, with early exit routing reducing average processing latency. Improves throughput and reduces average latency compared to naive sequential pipelines, with graceful degradation when the ML stage is unavailable.

    \item \textbf{Intelligent Evolution Controller} [RQ3]: A multi-strategy self-adaptation system that monitors quality meta-streams via drift detection and selects from multiple adaptation strategies based on observed drift characteristics. Resolves the majority of drift events through lightweight incremental updates without full retraining.

    \item \textbf{Online Autoencoder with Memory Module} [RQ4, RQ5]: A memory-augmented anomaly detector that learns multivariate normal patterns from ultra-clean training data and scores anomalies via reconstruction error, with continuous memory-based learning for streaming data without catastrophic forgetting. Achieves strong detection performance on multivariate anomalies.
\end{enumerate}

These contributions are not independent components but a tightly integrated system: the multi-dimensional threshold matrix provides context boundaries that the clusterer uses for assignment; the memory-augmented autoencoder scores anomalies within each context; the sequential four-stage pipeline with early exit routing meets throughput requirements through staged validation; and the drift controller orchestrates model updates without pipeline interruption.

% ================================================
\section{Thesis Content}
% ================================================

This thesis is organized into six chapters:

\begin{enumerate}
    \item \textbf{Chapter 1: Introduction} presents the motivation for streaming data quality monitoring, identifies three research challenges (context collapse, validation pipeline inefficiency, single-strategy drift adaptation), states the thesis objective, and highlights the four primary contributions.

    \item \textbf{Chapter 2: Background} reviews foundational concepts including streaming data processing (Apache Kafka, Apache Flink), data quality dimensions (DAMA-DMBOK), anomaly detection methods (Isolation Forest, Autoencoder), concept drift (ADWIN), and context-aware systems. This chapter establishes the technical infrastructure and conceptual vocabulary on which subsequent chapters build.

    \item \textbf{Chapter 3: Related Work} surveys four interconnected research areas: streaming data quality monitoring frameworks (Great Expectations, Soda Core, Confluent, Grab Coban, Stream DaQ), concept drift handling and ADWIN variants (ADWIN-U, Parallel ADWIN, ADWIN++), context-aware anomaly detection (ReAD, BeSTAD, ICAD, DTD), and memory-based streaming methods (Online Autoencoder with Memory, MStream). This chapter identifies the specific gaps that CA-DQStream addresses.

    \item \textbf{Chapter 4: CA-DQStream: Context-Aware Data Quality Streaming System} presents the complete system design including the approach direction, exploratory data analysis on the NYC Taxi dataset, the four-layer processing pipeline (Baseline Validation, Memory-Based Streaming Anomaly Detection, Meta-Aggregation, Intelligent Evolution Controller), feature engineering with 21-dimensional feature vectors, and fault tolerance mechanisms.

    \item \textbf{Chapter 5: Experiments and Evaluation} validates the four contributions through comprehensive experiments organized by Research Question (RQ1--RQ5) on the NYC Yellow Taxi dataset (72M records across 26 months: 14 months real TLC data + 12 months simulated data), demonstrating traceability between proposed innovations and empirical validation.

    \item \textbf{Chapter 6: Conclusion and Future Work} summarizes contributions, acknowledges limitations (single-domain evaluation, synthetic ground truth), and outlines promising directions for future research.
\end{enumerate}
